\newlabel{tab:short_long_run}{{1}{2}{Summary of AI's Impact in the Short, Medium, and Long Run\relax }{table.caption.1}{}}
\newlabel{fig:task_division}{{1}{3}{Illustrative Example for Division of Labor\relax }{figure.caption.3}{}}
\newlabel{fig:steps_tasks}{{2}{10}{Bundling Steps into Tasks via AI Strategy\relax }{figure.caption.9}{}}
\newlabel{sec:model.steps}{{3.1}{10}{Steps}{subsection.3.1}{}}
\newlabel{def:manual_task}{{1}{10}{Manual Step}{definition.1}{}}
\newlabel{def:augmented_task}{{2}{10}{Augmented Step}{definition.2}{}}
\newlabel{sec:model.tasks}{{3.2}{11}{Tasks}{subsection.3.2}{}}
\newlabel{def:automated_task}{{3}{11}{Automated Step}{definition.3}{}}
\newlabel{def:ai_chain}{{4}{11}{AI Chain}{definition.4}{}}
\newlabel{tab:model_parameters}{{2}{12}{Notation Summary\relax }{table.caption.10}{}}
\newlabel{sec:jobs_wages}{{3.3}{12}{Production}{subsection.3.3}{}}
\newlabel{eq:wage}{{1}{13}{Production}{equation.3.1}{}}
\newlabel{sec:model.jobs}{{3.4}{13}{Firm's Cost Minimization Problem}{subsection.3.4}{}}
\newlabel{foot:wage}{{9}{13}{}{Hfootnote.14}{}}
\newlabel{sec:model.shortrun}{{3.4.1}{14}{Short Run: Fixed Job Boundaries and Wages, Adjustable Execution Times}{subsubsection.3.4.1}{}}
\newlabel{eq:totalcost_shortterm}{{2}{14}{Short Run: Fixed Job Boundaries and Wages, Adjustable Execution Times}{equation.3.2}{}}
\newlabel{prop:shortrun_dp}{{1}{14}{Short-run AI Deployment}{proposition.1}{}}
\newlabel{sec:model.mediumrun}{{3.4.2}{15}{Medium Run: Fixed Job Boundaries, Adjustable Wages and Execution Times}{subsubsection.3.4.2}{}}
\newlabel{eq:totalcost_mediumrun}{{3}{15}{Medium Run: Fixed Job Boundaries, Adjustable Wages and Execution Times}{equation.3.3}{}}
\newlabel{prop:mediumrun_dp}{{2}{15}{Medium-run AI Deployment}{proposition.2}{}}
\newlabel{sec:model.longrun}{{3.4.3}{16}{Long Run: Adjustable Job Boundaries, Wages, and Execution Times}{subsubsection.3.4.3}{}}
\newlabel{sec:comparative_advantage}{{4.1}{16}{Overturning of Comparative Advantage due to Chaining}{subsection.4.1}{}}
\newlabel{eq:ca_test}{{4}{17}{Overturning of Comparative Advantage due to Chaining}{equation.4.4}{}}
\newlabel{ex:ca_predecessor}{{1}{18}{}{example.1}{}}
\newlabel{sec:fragmentation}{{4.2}{20}{Job-level AI Exposure and the Fragmentation Index}{subsection.4.2}{}}
\newlabel{ex:fragmentation_interleaved}{{2}{20}{}{example.2}{}}
\newlabel{ex:fragmentation_clustered}{{3}{21}{}{example.3}{}}
\newlabel{eq:fragmentation}{{5}{23}{Job-level AI Exposure and the Fragmentation Index}{equation.4.5}{}}
\newlabel{prop:fragmentation}{{3}{23}{}{proposition.3}{}}
\newlabel{sec:nonlinear}{{4.3}{23}{Non-Monotonicity of Gains from AI Improvements}{subsection.4.3}{}}
\newlabel{ex:nonmonotone}{{4}{24}{}{example.4}{}}
\newlabel{fig:example5_nonmonotone}{{3}{25}{Optimal Cost and the Marginal Benefit of AI Quality in Example~\ref {ex:nonmonotone}\relax }{figure.caption.11}{}}
\newlabel{fig:example5_costs}{{3(a)}{25}{Total cost by AI strategy\relax }{figure.caption.11}{}}
\newlabel{sub@fig:example5_costs}{{(a)}{25}{Total cost by AI strategy\relax }{figure.caption.11}{}}
\newlabel{fig:example5_marginalBenefit}{{3(b)}{25}{Marginal benefit of improving $\alpha $\relax }{figure.caption.11}{}}
\newlabel{sub@fig:example5_marginalBenefit}{{(b)}{25}{Marginal benefit of improving $\alpha $\relax }{figure.caption.11}{}}
\newlabel{sec:empirics}{{5}{26}{Empirical Evaluation}{section.5}{}}
\newlabel{sec:chainLength_prediction}{{5.1}{27}{Prediction \#1: Tendency to Have Runs of Consecutive AI-executed Tasks}{subsection.5.1}{}}
\newlabel{fig:aiChains_graphs_def1}{{4}{28}{Average Length of AI Chains: Actual versus Placebo Datasets\relax }{figure.caption.12}{}}
\newlabel{sec:fragmentation_prediction}{{5.2}{29}{Prediction \#2: Dispersion of AI-able Tasks Affects AI Execution Outcomes}{subsection.5.2}{}}
\newlabel{eq:fragmentation_index_regression}{{6}{29}{Prediction \#2: Dispersion of AI-able Tasks Affects AI Execution Outcomes}{equation.5.6}{}}
\newlabel{tab:fragmentation_index_regression_exposure}{{3}{31}{Role of AI-able Task Fragmentation in Determining AI Execution Outcomes\relax }{table.caption.13}{}}
\newlabel{sec:DWA_prediction}{{5.3}{31}{Prediction \#3: Adjacency to AI Tasks Increases Likelihood of AI Execution}{subsection.5.3}{}}
\newlabel{eq:DWA_regression_ai}{{7}{32}{Prediction \#3: Adjacency to AI Tasks Increases Likelihood of AI Execution}{equation.5.7}{}}
\newlabel{tab:DWA_regression_aiExecution_mainSample}{{4}{33}{Effect of Neighboring Tasks' AI Execution Status on Task's AI Execution\relax }{table.caption.14}{}}
\newlabel{sec:robustness}{{5.4}{34}{Robustness Tests}{subsection.5.4}{}}
\newlabel{fig:DWA_regression_aiExecution_mainSample}{{5}{35}{Effect of Neighboring Tasks' AI Execution Status on Task's AI Execution\relax }{figure.caption.15}{}}
\newlabel{fig:ame_nofe}{{5(a)}{35}{No Fixed Effects\relax }{figure.caption.15}{}}
\newlabel{sub@fig:ame_nofe}{{(a)}{35}{No Fixed Effects\relax }{figure.caption.15}{}}
\newlabel{fig:ame_major}{{5(b)}{35}{SOC Major Group Fixed Effects\relax }{figure.caption.15}{}}
\newlabel{sub@fig:ame_major}{{(b)}{35}{SOC Major Group Fixed Effects\relax }{figure.caption.15}{}}
\newlabel{fig:ame_minor}{{5(c)}{35}{SOC Minor Group Fixed Effects\relax }{figure.caption.15}{}}
\newlabel{sub@fig:ame_minor}{{(c)}{35}{SOC Minor Group Fixed Effects\relax }{figure.caption.15}{}}
\newlabel{fig:ame_dwa}{{5(d)}{35}{Detailed Work Activity Fixed Effects\relax }{figure.caption.15}{}}
\newlabel{sub@fig:ame_dwa}{{(d)}{35}{Detailed Work Activity Fixed Effects\relax }{figure.caption.15}{}}
\newlabel{app:omitted_proofs}{{A}{OA - 1}{Omitted Proofs}{appendix.A}{}}
\newlabel{lem:FI.upper.bound}{{OA-1}{OA - 3}{}{lemma.OA-1}{}}
\newlabel{ex:FI.tight}{{OA-1}{OA - 5}{}{example.1}{}}
\newlabel{lem:FI.lower.bound.4}{{OA-2}{OA - 5}{}{lemma.OA-2}{}}
\newlabel{ex:FI.lower.gap}{{OA-2}{OA - 7}{}{example.2}{}}
\newlabel{lem:FI.lower.bound.8}{{OA-3}{OA - 7}{}{lemma.OA-3}{}}
\newlabel{ex:FI.necessity}{{OA-3}{OA - 8}{}{example.3}{}}
\newlabel{app:nonmonotone}{{A.3}{OA - 9}{Non-Monotone Gains from AI Improvements: Formal Statement and Proof}{subsection.A.3}{}}
\newlabel{eq:strategy_cost}{{OA.1}{OA - 9}{Non-Monotone Gains from AI Improvements: Formal Statement and Proof}{equation.A.1}{}}
\newlabel{lem:reorg_threshold}{{OA-4}{OA - 9}{Non-monotone Marginal Returns to AI Quality Improvements}{lemma.OA-4}{}}
\newlabel{fig:hierarchical_production}{{OA-1}{OA - 11}{Illustrative Example of a Firm's Production in the Long Run\relax }{figure.caption.16}{}}
\newlabel{app:jobdesign}{{B}{OA - 11}{Job Design and the Firm's Long Run Problem}{appendix.B}{}}
\newlabel{sec:app.jobs}{{B.1}{OA - 11}{Jobs, Wages, and Hand-off Costs}{subsection.B.1}{}}
\newlabel{eq:wage_jobdesign}{{OA.2}{OA - 12}{Jobs, Wages, and Hand-off Costs}{equation.B.2}{}}
\newlabel{eq:wagebill}{{OA.3}{OA - 12}{Jobs, Wages, and Hand-off Costs}{equation.B.3}{}}
\newlabel{fig:intro_task_illustration}{{OA-2}{OA - 13}{Illustration of the Trade-off between Worker Specialization and Coordination in Task Assignment\relax }{figure.caption.17}{}}
\newlabel{tab:jobdesign_parameters}{{OA-1}{OA - 13}{Additional Notation for Job Design\relax }{table.caption.18}{}}
\newlabel{app:jobdesign.longrun}{{B.2}{OA - 13}{Firm's Long-Run Optimization Problem}{subsection.B.2}{}}
\newlabel{eq:totalcost_with_handoff}{{OA.4}{OA - 14}{Firm's Long-Run Optimization Problem}{equation.B.4}{}}
\newlabel{prop:totalcost_optimization_dp}{{OA-1}{OA - 14}{Long-run AI Deployment and Job Design}{proposition.OA-1}{}}
\newlabel{sec:model_discussion}{{B.3}{OA - 16}{Hand-off Costs and the Limits of Worker Specialization}{subsection.B.3}{}}
\newlabel{tab:job_design}{{OA-2}{OA - 17}{Role of Hand-off Costs in Organizational Structure\relax }{table.caption.19}{}}
\newlabel{fig:job_design}{{OA-3}{OA - 18}{Geometric Interpretation of Job Designs\relax }{figure.caption.20}{}}
\newlabel{sec:aggregation}{{C}{OA - 20}{Macro-level Production Function}{appendix.C}{}}
\newlabel{eq:new_wage}{{OA.5}{OA - 21}{Macro-level Production Function}{equation.C.5}{}}
\newlabel{eq:skill_adjusted_time}{{C}{OA - 21}{Macro-level Production Function}{equation.C.5}{}}
\newlabel{sec:agg_within}{{C.1}{OA - 22}{Within-Firm Aggregation: Leontief to Leontief}{subsection.C.1}{}}
\newlabel{eq:task_level_prod}{{OA.6}{OA - 23}{Within-Firm Aggregation: Leontief to Leontief}{equation.C.6}{}}
\newlabel{eq:task_subst_rate_indep}{{OA.7}{OA - 23}{Within-Firm Aggregation: Leontief to Leontief}{equation.C.7}{}}
\newlabel{eq:micro_agg_prod}{{OA.8}{OA - 23}{Within-Firm Aggregation: Leontief to Leontief}{equation.C.8}{}}
\newlabel{eq:def_t_AI}{{OA.9}{OA - 24}{Within-Firm Aggregation: Leontief to Leontief}{equation.C.9}{}}
\newlabel{eq:def_t_H}{{OA.10}{OA - 24}{Within-Firm Aggregation: Leontief to Leontief}{equation.C.10}{}}
\newlabel{eq:alpha_bar}{{OA.11}{OA - 24}{Within-Firm Aggregation: Leontief to Leontief}{equation.C.11}{}}
\newlabel{eq:micro_fixed_proportions}{{OA.12}{OA - 24}{Within-Firm Aggregation: Leontief to Leontief}{equation.C.12}{}}
\newlabel{eq:macro_agg_prod}{{OA.13}{OA - 26}{Cross-Firm Aggregation: Leontief to CES}{equation.C.13}{}}
\newlabel{eq:alpha_threshold}{{OA.14}{OA - 27}{Cross-Firm Aggregation: Leontief to CES}{equation.C.14}{}}
\newlabel{eq:micro_X}{{OA.15}{OA - 27}{Cross-Firm Aggregation: Leontief to CES}{equation.C.15}{}}
\newlabel{eq:micro_M}{{OA.16}{OA - 27}{Cross-Firm Aggregation: Leontief to CES}{equation.C.16}{}}
\newlabel{eq:micro_H}{{OA.17}{OA - 27}{Cross-Firm Aggregation: Leontief to CES}{equation.C.17}{}}
\newlabel{eq:macro_agg_prod_with_micro_aggregates}{{OA.18}{OA - 27}{Cross-Firm Aggregation: Leontief to CES}{equation.C.18}{}}
\newlabel{eq:phi}{{OA.19}{OA - 27}{Cross-Firm Aggregation: Leontief to CES}{equation.C.19}{}}
\newlabel{sec:aggregationDist}{{C.3}{OA - 28}{Derivation of the Effective AI Quality Heterogeneity Distribution}{subsection.C.3}{}}
\newlabel{eq:macro_agg_prod_with_hetDist}{{OA.20}{OA - 28}{Derivation of the Effective AI Quality Heterogeneity Distribution}{equation.C.20}{}}
\newlabel{eq:gamma}{{OA.21}{OA - 28}{Derivation of the Effective AI Quality Heterogeneity Distribution}{equation.C.21}{}}
\newlabel{eq:psi}{{OA.22}{OA - 28}{Derivation of the Effective AI Quality Heterogeneity Distribution}{equation.C.22}{}}
\newlabel{eq:macro_agg_pro_with_gamma}{{OA.23}{OA - 28}{Derivation of the Effective AI Quality Heterogeneity Distribution}{equation.C.23}{}}
\newlabel{eq:psi_solved}{{OA.24}{OA - 29}{Derivation of the Effective AI Quality Heterogeneity Distribution}{equation.C.24}{}}
\newlabel{app:data_prep}{{D}{OA - 30}{Construction Details of the Main Sample}{appendix.D}{}}
\newlabel{fig:task_execution_dist}{{OA-4}{OA - 31}{Distribution of Modes of Task Execution in the Dataset\relax }{figure.caption.21}{}}
\newlabel{fig:occupation_ai_share}{{OA-5}{OA - 32}{Distribution of Share of Occupation Tasks Exposed to and Executed by AI\relax }{figure.caption.22}{}}
\newlabel{fig:computer_programmers}{{OA-6}{OA - 32}{Task Sequence of Computer Programmers Occupation\relax }{figure.caption.23}{}}
\newlabel{fig:public_relations}{{OA-7}{OA - 33}{Task Sequence of Public Relations Specialists Occupation\relax }{figure.caption.24}{}}
\newlabel{fig:electronic_equipments}{{OA-8}{OA - 34}{Task Sequence of Electronic Equipment Installers and Repairers, Motor Vehicles\relax }{figure.caption.25}{}}
\newlabel{tab:example_anthropic_AI_tasks}{{OA-3}{OA - 36}{Example O*NET Tasks Associated with Claude Conversations, and Calculated Augmentation vs.\ Automation Outcomes\relax }{table.caption.26}{}}
\newlabel{tab:fragmentation_index_regression_execution}{{OA-4}{OA - 37}{Role of AI-able Task Fragmentation in Determining AI Execution Outcomes (Execution-Based Empirical Fragmentation Measure)\relax }{table.caption.27}{}}
\newlabel{app:additional_robustness}{{E}{OA - 37}{Additional Robustness Tests for the Empirical Section}{appendix.E}{}}
\newlabel{app:additional_robustness_pred2}{{E.1}{OA - 37}{Additional Tests for Prediction \#2}{subsection.E.1}{}}
\newlabel{app:additional_robustness_pred3}{{E.2}{OA - 38}{Additional Tests for Prediction \#3}{subsection.E.2}{}}
\newlabel{tab:DWA_regression_aiExecution_GPTsample}{{OA-5}{OA - 39}{Effect of Neighboring Tasks' AI Execution Status on Task's AI Execution Likelihood (GPT-filtered Sample)\relax }{table.caption.28}{}}
\newlabel{tab:DWA_regression_aiAutomation_mainSample}{{OA-6}{OA - 40}{Effect of Neighboring Tasks' AI Execution Status on Task's AI Automation Likelihood\relax }{table.caption.29}{}}
\newlabel{tab:DWA_regression_aiAutomation_GPTsample}{{OA-7}{OA - 41}{Effect of Neighboring Tasks' AI Execution Status on Task's AI Automation Likelihood (GPT-filtered Sample)\relax }{table.caption.30}{}}
\newlabel{fig:DWA_regression_aiExecution_GPTsample}{{OA-9}{OA - 42}{Effect of Neighboring Tasks' AI Execution Status on Task's AI Execution Likelihood (GPT-filtered Sample)\relax }{figure.caption.31}{}}
\newlabel{fig:DWA_regression_aiAutomation_mainSample}{{OA-10}{OA - 43}{Effect of Neighboring Tasks' AI Execution Status on Task's AI Automation Likelihood\relax }{figure.caption.32}{}}
\newlabel{fig:DWA_regression_aiAutomation_GPTsample}{{OA-11}{OA - 44}{Effect of Neighboring Tasks' AI Execution Status on Task's AI Automation Likelihood (GPT-filtered Sample)\relax }{figure.caption.33}{}}
\newlabel{tab:DWA_tasks}{{OA-8}{OA - 45}{Example DWAs with Tasks in Multiple Occupations and the Tasks Surviving the ``Finding Similar Tasks'' Procedure\relax }{table.caption.34}{}}
\newlabel{app:prompts}{{F}{OA - 46}{GPT-5-mini Prompts}{appendix.F}{}}
\newlabel{app:gptPrompts_robustness}{{G}{OA - 48}{Robustness to Alternative GPT Prompts}{appendix.G}{}}
\newlabel{fig:kendall_tau_dist}{{OA-12}{OA - 49}{Kendall's $\tau $ Distribution Across GPT Prompt Task Orderings\relax }{figure.caption.35}{}}
\newlabel{fig:ai_chains_length_count_prompt_robustness}{{OA-13}{OA - 50}{Robustness of AI Chain Length to GPT Prompts\relax }{figure.caption.36}{}}
\newlabel{app:gptPrompts_robustness_pred1}{{G.1}{OA - 50}{Robustness of Prediction \#1 Results to Alternative GPT Prompts}{subsection.G.1}{}}
\newlabel{fig:efi_regression_exposure_robustness}{{OA-14}{OA - 51}{Robustness of Table~\ref {tab:fragmentation_index_regression_exposure} Estimates to GPT Prompts\relax }{figure.caption.37}{}}
\newlabel{app:gptPrompts_robustness_pred2}{{G.2}{OA - 51}{Robustness of Prediction \#2 Results to Alternative GPT Prompts}{subsection.G.2}{}}
\newlabel{app:gptPrompts_robustness_pred3}{{G.3}{OA - 51}{Robustness of Prediction \#3 Results to Alternative GPT Prompts}{subsection.G.3}{}}
\newlabel{fig:DWA_regression_aiExecution_mainSample_robustness}{{OA-15}{OA - 53}{Robustness of Table~\ref {tab:DWA_regression_aiExecution_mainSample} Estimates to GPT Prompts\relax }{figure.caption.38}{}}
\newlabel{app:frequency_robustness}{{H}{OA - 54}{Robustness to Frequently-Executed Tasks Sample Restriction}{appendix.H}{}}
\newlabel{app:frequency_robustness_construction}{{H.1}{OA - 55}{Constructing the Frequent-Task Samples}{subsection.H.1}{}}
\newlabel{app:frequency_robustness_chains}{{H.2}{OA - 56}{Prediction \#1: The Average AI Chain Length}{subsection.H.2}{}}
\newlabel{fig:frequency_pruning_example}{{OA-16}{OA - 57}{Frequency Pruning of a Single Occupation's Workflow: Allergists and Immunologists\relax }{figure.caption.39}{}}
\newlabel{fig:chainlength_frequency}{{OA-17}{OA - 58}{Average AI Chain Length Versus Its Task Position-Reshuffle Null Across Frequency Cuts\relax }{figure.caption.40}{}}
\newlabel{app:frequency_robustness_frag}{{H.3}{OA - 58}{Prediction \#2: The Empirical Fragmentation Index}{subsection.H.3}{}}
\newlabel{fig:frag_frequency}{{OA-18}{OA - 59}{Empirical Fragmentation Index and AI Execution Across Frequency Cuts\relax }{figure.caption.41}{}}
\newlabel{app:frequency_robustness_neighbor}{{H.4}{OA - 59}{Prediction \#3: The Immediate-Neighbor Effect}{subsection.H.4}{}}
\newlabel{fig:neighbor_frequency_heatmap}{{OA-19}{OA - 61}{Neighbor AI Execution and Focal-Task AI Execution Across Frequency Cuts\relax }{figure.caption.42}{}}
\newlabel{fig:neighbor_placebo_forest}{{OA-20}{OA - 62}{Immediate-Neighbor Effects: Observed Marginal Effects Versus Their Position-Reshuffle Nulls Across Frequency Cuts\relax }{figure.caption.43}{}}
\newlabel{app:external_validation}{{I}{OA - 63}{External Validation of the Sequencing Results}{appendix.I}{}}
\newlabel{tab:apqc_process_examples}{{OA-11}{OA - 64}{Example Process Groups From Two Industry Frameworks\relax }{table.caption.44}{}}
\newlabel{app:apqc_validation}{{I.1}{OA - 64}{Concern \#1: LLM Recovering Practitioner-Documented Task Orderings}{subsection.I.1}{}}
\newlabel{app:apqc_validation_caveat}{{I.1}{OA - 65}{Practitioner orderings are not immutable ground truth}{paragraph*.45}{}}
\newlabel{app:apqc_validation_design}{{I.1}{OA - 65}{Design}{paragraph*.46}{}}
\newlabel{tab:apqc_pcf_ordering_validation}{{OA-12}{OA - 66}{Overlap Between GPT-Generated and Practitioner-Documented Process Orderings\relax }{table.caption.48}{}}
\newlabel{app:apqc_validation_results}{{I.1}{OA - 66}{Results}{paragraph*.47}{}}
\newlabel{fig:apqc_tau_by_size}{{OA-21}{OA - 67}{Overlap between GPT-orderings and Original APQC-PCF Order by Branch Length\relax }{figure.caption.49}{}}
\newlabel{fig:apqc_tau_size_main}{{OA-21(a)}{OA - 67}{Main prompt\relax }{figure.caption.49}{}}
\newlabel{sub@fig:apqc_tau_size_main}{{(a)}{OA - 67}{Main prompt\relax }{figure.caption.49}{}}
\newlabel{fig:apqc_tau_size_alt}{{OA-21(b)}{OA - 67}{Average of ten alternative prompts\relax }{figure.caption.49}{}}
\newlabel{sub@fig:apqc_tau_size_alt}{{(b)}{OA - 67}{Average of ten alternative prompts\relax }{figure.caption.49}{}}
\newlabel{app:apqc_validation_tail}{{I.1}{OA - 67}{Reading the left tail}{paragraph*.50}{}}
\newlabel{app:apqc_fragmentation_top}{{I.2}{OA - 68}{Concern \#2: Re-Estimating the Model Predictions on Documented Sequences}{subsection.I.2}{}}
\newlabel{tab:apqc_fragmentation_index_regression}{{OA-13}{OA - 69}{Fragmentation and AI Execution on APQC's Documented Process Sequences\relax }{table.caption.54}{}}
\newlabel{app:external_validation_discussion}{{I.3}{OA - 70}{Summary and Takeaways}{subsection.I.3}{}}
